\documentclass[11pt]{article}

\usepackage[left=0.82in,right=0.82in,top=0.72in,bottom=0.78in]{geometry}
\usepackage{amsmath}
\usepackage{booktabs}
\usepackage{enumitem}
\usepackage{fancyhdr}
\usepackage{graphicx}
\usepackage{microtype}
\usepackage{xcolor}
\usepackage[hidelinks]{hyperref}

\definecolor{programblue}{HTML}{183B56}
\definecolor{programteal}{HTML}{176B6B}
\definecolor{softblue}{HTML}{EDF4F8}
\definecolor{softgray}{HTML}{F3F4F5}
\definecolor{statusgreen}{HTML}{1F6A43}
\definecolor{statusamber}{HTML}{8A5B00}

\hypersetup{
  colorlinks=true,
  linkcolor=programblue,
  urlcolor=programteal
}
\setlist[itemize]{leftmargin=1.25em,itemsep=2pt,topsep=4pt}
\setlist[enumerate]{leftmargin=1.45em,itemsep=3pt,topsep=4pt}
\setlength{\parindent}{0pt}
\setlength{\parskip}{5pt}
\renewcommand{\arraystretch}{1.14}

\pagestyle{fancy}
\fancyhf{}
\lhead{\small First-principles open dynamics}
\rhead{\small Repository overview}
\cfoot{\small \thepage}
\renewcommand{\headrulewidth}{0.3pt}

\newcommand{\implemented}{\textcolor{statusgreen}{\textbf{Implemented and internally verified}}}
\newcommand{\fixturevalidated}{\textcolor{statusamber}{\textbf{Fixture-validated}}}
\newcommand{\proposed}{\textcolor{programblue}{\textbf{Proposed}}}
\newcommand{\independent}{\textcolor{statusamber}{\textbf{Independent and unreviewed}}}

\begin{document}

\begin{center}
  {\Large\bfseries\color{programblue}
  Repository Alignment with the Ping Group's\\[2pt]
  First-Principles Open-Dynamics Program}

  \vspace{6pt}
  {\large Scientific direction, implemented first slice, and collaboration gates}

  \vspace{5pt}
  {\small August 2026}
\end{center}

\vspace{4pt}
\noindent\colorbox{softblue}{%
  \parbox{0.965\linewidth}{%
  \textbf{Decision.} I have reorganized the repository around a
  source-anchored benchmark for equal-time non-Markovian electron--phonon
  dynamics and a future first-principles interchange layer. Quantum algorithms
  now enter as downstream, falsifiable selected-mode benchmarks: ADAPT-VQE for
  stationary state preparation, adaptive McLachlan for explicit selected-mode
  time evolution, and QSE for spectra or response subspaces. They are not
  presented as replacements for first-principles density-matrix dynamics.}}

\section*{1. Scientific center: equal-time non-Markovian electron--phonon dynamics}

The immediate scientific anchor is the first-principles density-matrix program
described by Simoni, Riva, and Ping: electronic and phonon structure;
interaction matrix elements; and Markovian or non-Markovian density-matrix
propagation.\footnote{J.~Simoni, G.~Riva, and Y.~Ping,
\href{https://arxiv.org/html/2504.17936v2}{\emph{First-principles open quantum
dynamics for solids based on density-matrix formalism}} (2025).}
The new equal-time formulation evolves the electronic one-body density matrix
\(\rho\), coherent phonon amplitudes \(B\), normal and anomalous phonon
fluctuations \(N\) and \(A\), and connected electron--phonon correlations
\(C\).\footnote{G.~Riva, J.~Simoni, and Y.~Ping,
\href{https://arxiv.org/html/2606.22233v1}{\emph{Open-quantum-system theory of
non-Markovian electron--phonon dynamics}} (2026).}

The repository already contained two useful ends of a validation hierarchy:
an explicit truncated Holstein-dimer propagator and an independent
31-real-coordinate implementation of the complete two-site matrix variables.
The first new integration layer now places solver outputs in the common object
\[
\mathcal T=\{t_j,\rho(t_j);\ B(t_j),N(t_j),A(t_j),C(t_j);\
\text{method, units, conventions, provenance}\}.
\]
Only available moments are present. Missing collision or fluctuation physics is
not represented by zeros.

\textbf{Claim boundary.} The public v1 arXiv archive does not contain the
supplement identified in the article as giving the spin-reduced dimer
equations, and it contains no numerical dataset. I therefore describe this as
a \emph{source-anchored independent invariant benchmark}, not a reproduction
of the authors' implementation or curves. The representability and closure
diagnoses are my independent work; Gabriele, Jacopo, and Professor Ping have
not reviewed or endorsed them.

\section*{2. Implemented first slice}

\begin{tabular}{p{0.23\linewidth}p{0.47\linewidth}p{0.20\linewidth}}
\toprule
Object & Concrete behavior & Status \\
\midrule
Reduced trajectory & Required electronic 1-RDM and basis; optional coherent,
normal, anomalous, and connected moments; explicit units and method identity.
& \implemented \\
\addlinespace
Reference isolation & Exact trajectories are labeled reporting-only and expose
neither wavefunctions nor Hamiltonians to adaptive decisions.
& \implemented \\
\addlinespace
Diagnostics & Structural errors raise; Hermiticity, trace, eigenvalue, and
phonon-symmetry failures are reported without clipping, projection, or sample
deletion.
& \implemented \\
\addlinespace
Dimer adapter & Field-wise reuse of the existing exact, coherent-only, and
31-coordinate implementations on one sampled grid and common initial reduced
moments.
& \implemented \\
\addlinespace
Protocol record & Paper-stated parameters are stored separately from the
repository's cutoff, integration, initialization, tolerance, and Fourier
choices.
& \fixturevalidated \\
\bottomrule
\end{tabular}

The first implementation is covered by 12 focused tests. An additional 19
unchanged tests verify the underlying matrix-coordinate parity and exact driven
reference. This establishes an internally verified interoperability seam; it
does not establish external or real-material validation.

\begin{center}
\fbox{\begin{minipage}{0.93\linewidth}
\centering
\textbf{Implemented comparison path}\\[4pt]
truncated exact reference (offline only)\quad|\quad coherent-only dynamics
\quad|\quad independent equal-time dynamics\\[3pt]
\(\Downarrow\)\\[-1pt]
common reduced moments, polarization, declared Hann spectrum, and physicality
diagnostics
\end{minipage}}
\end{center}

\section*{3. How the existing projects now serve this program}

The earlier projects become capabilities within the open-dynamics program
instead of separate claims competing for attention.

\begin{tabular}{p{0.17\linewidth}p{0.28\linewidth}p{0.44\linewidth}}
\toprule
Capability & Scientific role & Direct group-facing use \\
\midrule
Resource-aware ADAPT-VQE & Stationary correlated-state construction by
geometry- and cost-aware candidate-position selection. & Prepare a compact
initial state for a selected electronic-plus-phonon active problem. It is not a
dynamics method. \\
\addlinespace
Adaptive McLachlan & Real-time maintenance of a variational manifold with
append, prune, and exchange updates. & Propagate an explicit selected-mode
supersystem and emit the same reduced moments as the classical benchmark. \\
\addlinespace
QSE and response & Excited or response subspaces and spectral observables. &
Supply response diagnostics or selected excited-state information; do not
identify QSE automatically with full real-time propagation. \\
\addlinespace
Molecular vibronics & Psi4-derived water fixture and a matched PySCF
electronic/harmonic cross-check. & Test the active-problem interface away from
the analytic dimer. Electron--phonon derivatives and full dynamics still need
separate validation. \\
\addlinespace
Representability analysis & Joint electronic--bosonic Gram diagnostics and
offline closure comparisons. & Diagnose when an independently implemented
reduced trajectory leaves necessary physicality cones; no claim about a defect
in the published theory. \\
\bottomrule
\end{tabular}

\section*{4. The quantum-computing question}

The near-term quantum question is not whether a circuit can replace the full
FPDMD workflow. It is whether a selected explicit electron--phonon subsystem
can provide a controlled benchmark or closure diagnostic for non-Markovian
dynamics. The proposed data flow is
\[
\text{material data}
\longrightarrow \text{selected electrons and modes}
\longrightarrow
\begin{cases}
\text{ADAPT-VQE initial state},\\
\text{adaptive McLachlan trajectory},\\
\text{QSE response information},
\end{cases}
\longrightarrow \mathcal T.
\]
The output is not merely an energy. It is the common set of measured
electronic, phononic, and connected moments, along with circuit depth,
measurement work, phonon truncation, and uncertainty.

The classical baselines are explicit truncated evolution, coherent-only
dynamics, and the equal-time non-Markovian model. Exact reference data remain
unavailable to operator selection, generator admission, pruning, stopping, and
regularization. A quantum lane should advance beyond exact-solvable fixtures
only if shared moments agree within joint numerical and sampling uncertainty,
improve on the coherent-only baseline at preregistered held-out points, and
survive timestep and phonon-cutoff refinement. Passing this gate supports a
larger benchmark; it is not evidence of quantum advantage.

\section*{5. First-principles and group-native interfaces}

\subsection*{5.1 Portable material bundle}

The group's public FPDMD workflow is currently native to JDFTx and a modified
FeynWann initializer that produces \texttt{ldbd\_data}.\footnote{
\href{https://github.com/Ping-Group-UCSC/FPDMD_development}{Public FPDMD code}
and the
\href{https://github.com/Ping-Group-UCSC/Tutorials/blob/master/DMD-workflows/GaAs-DMD-example.md}{public GaAs workflow}.}
The legacy binary directory is not a safe external contract: public code
inspection shows native binary layouts whose shapes and interpretation depend
on positional metadata and runtime choices. The first group-native interface
should therefore be a producer-side portable FeynWann export, not a guessed
binary parser.

The proposed neutral bundle carries electronic energies and occupations,
phonon frequencies and eigenvectors, full complex electron--phonon matrix
elements, mesh weights, \(k+q\) maps, units, gauge/basis conventions, symmetry
expansion, zero-point and finite-grid normalization, active windows, and
software/input hashes. Quantum ESPRESSO, DFPT, Wannier90, and EPW can become
alternate producers of the same contract. Full complex couplings require a
version-pinned exporter; printed coupling magnitudes cannot support coherent
phase-sensitive dynamics.

\textbf{Go gate:} implement this bundle against a genuine small producer-side
fixture with complete provenance. Until then the interface remains proposed
and is not called FPDMD-, FeynWann-, QE-, or EPW-compatible.

\subsection*{5.2 INQ and coherent real-time observables}

Jacopo and collaborators' spin-noncollinear real-time TDDFT implementation in
INQ supplies a natural coherent baseline and observable source.\footnote{
J.~Simoni et al.,
\href{https://arxiv.org/abs/2506.21908}{\emph{Spin non-Collinear Real-Time
Time-Dependent Density-Functional Theory and Implementation in the Modern
GPU-Accelerated INQ Code}} (2025).}
A future adapter would ingest declared time grids, fields, current,
polarization, spin, or magnetization traces into the common trajectory
comparison. It would not be described as an open-system TDDFT or as the
equal-time collision theory unless an environmental treatment is actually
present.

\textbf{Go gate:} add the adapter from a versioned, producer-generated trace
with basis, field, unit, and observable conventions.

\section*{6. Additional public workstreams: narrow future seams}

These are possible later extensions of the same contracts, not commitments or
claims of collaboration.

\begin{tabular}{p{0.22\linewidth}p{0.25\linewidth}p{0.42\linewidth}}
\toprule
Public workstream & Publicly associated researchers & Narrow repository seam \\
\midrule
Exciton--phonon and radiative dynamics & Gabriele Riva, Jacopo Simoni, Chunhao
Guo (\href{https://arxiv.org/abs/2504.18071}{source}) &
Exciton-basis energies, dipoles, and exciton--phonon couplings feeding the same
selected-mode and reduced-trajectory contracts. \\
\addlinespace
Magnon--phonon coupling & Wuzhang Fang, Jacopo Simoni, Yuan Ping
(\href{https://arxiv.org/abs/2501.14239}{source}) & Magnon bands,
phonon modes, and exchange-derivative couplings for a small avoided-crossing
fixture before any open-magnon claim. \\
\addlinespace
Real-time photoresponse and chirality & Jacopo Simoni, Junting Yu, Shihao Tu,
and public coauthors
(\href{https://arxiv.org/abs/2601.01059}{photoresponse};
\href{https://arxiv.org/abs/2504.11596}{chirality}) & Current,
polarization, harmonic, and angular-momentum traces under explicit field and
Fourier conventions. \\
\addlinespace
Spin defects and decoherence & Ping-group defect and transport workstreams &
Ingest computed transition or relaxation data into a few-mode coherence model;
do not recreate the group's specialized first-principles tools. \\
\bottomrule
\end{tabular}

\section*{7. Near-term milestones and concrete discussion requests}

\begin{enumerate}
  \item \textbf{Completed now:} common reduced trajectories, source/repository
  provenance separation, exact-reference isolation, the three-level dimer
  adapter, common observables, and focused tests.
  \item \textbf{Next validated fixture:} obtain one small, sanitized portable
  export from FeynWann, or a pinned QE/EPW export with full complex couplings,
  inputs, and hashes.
  \item \textbf{First active problem:} agree on a small material, band window,
  and one or two physically important phonon modes; establish convergence
  against the unselected modes before attaching a residual bath.
  \item \textbf{Quantum falsification study:} compile the same active problem
  into ADAPT-VQE and adaptive McLachlan, emit common moments, and compare
  offline under predeclared error and resource metrics.
\end{enumerate}

The most useful immediate discussion with Professor Ping and Jacopo would be:
(i) whether the common dimer benchmark is a useful validation harness;
(ii) which producer and smallest dataset should define the first portable
material fixture; (iii) which selected phonon modes and observables would be
scientifically meaningful; and (iv) whether the missing spin-reduced
supplement or a reference trajectory can be made available for a true parity
test.

\section*{8. Status language for any future summary}

\begingroup
\small
\setlength{\parskip}{2pt}
\implemented: supported by named internal tests; no external validation is
implied.\par
\fixturevalidated: demonstrated only on the declared fixture and conventions.\par
\independent: not reviewed, validated, or endorsed by the Ping group.\par
\proposed: specification or experiment only.\par
\textbf{Collaborator-validated}: a supplied artifact passes a predeclared test,
and the collaborator approves that specific compatibility description.
\endgroup

\end{document}
